\section{The lattice intersection and its cubic symmetries}
\label{sec:B}

We use positive definite lattices with roots of squared length two. Thus
\(\det L\) is the determinant of a Gram matrix, and
\(\Theta_L(\tau)=\sum_{x\in L}q^{(x,x)/2}\), where
\(q=e^{2\pi i\tau}\). The quadratic discriminant form is
\(Q_L(x+L)=(x,x)/2\pmod{\Z}\). All matrices in our certificates
are rational or integral. The two different quadratic forms sometimes
called the discriminant form, \((x,x)\pmod{2\Z}\) and \(Q_L\),
therefore differ by the displayed factor two.

\subsection{An independent construction and the eight extensions}

Put
\[
 A_5=\{x\in\Z^6:\ \textstyle\sum x_i=0\},\qquad
 D_4=\{y\in\Z^4:\ \textstyle\sum y_i\in2\Z\},\qquad
 R=A_5^4\perp D_4.
\]
The scalar product is the ordinary Euclidean product on the ambient
\(28\)-space. Write \(\lambda=(5,-1,-1,-1,-1,-1)/6\), and
\(\omega_j=\sum_{i=1}^j e_i-(j/6)\sum_{i=1}^6e_i\).
The five glue vectors, with their first four entries denoting multiples
of \(\lambda\), are
\begin{equation}\label{eq:B-glue}
\begin{split}
 &(3,3,3,3;0),\quad(3,3,0,0;(1,0,0,0)),\\
 &(3,0,3,0;(1/2,1/2,1/2,1/2)),\\
 &(2,2,2,0;0),\quad(2,4,0,2;0).
\end{split}
\end{equation}
We define \(N\) to be their span together with \(R\), and put
\[
 \rho=((5,3,1,-1,-3,-5)/2)^4\,\oplus\,(3,2,1,0),
 \qquad \alpha=(e_1-e_2,0,0,0;0).
\]
Then \((\rho,\alpha)=1\), \(\rho^2=84\), and \(\rho\in N\).
For \(d\mid6\) define, without any choice of a lift modulo \(d^2\),
\begin{equation}\label{eq:B-neighbor}
 K_d=\{x\in N:(\rho,x)\in d\Z\},\qquad
 N_d=K_d+\Z\bigl(\rho/d-(6/d)\alpha\bigr).
\end{equation}
Let \(r\) cyclically shift the six coordinates in every \(A_5\) factor
by two places and fix the \(D_4\) factor. Set
\(\Lambda_i=r^iN_6\), \(i\in\Z/3\Z\), and
\(J=\Lambda_0\cap\Lambda_1\).

\begin{statement}[B.1 (Proved): the intersection and its discriminant form]
The lattice \(N\) is even unimodular, with root system \(A_5^4D_4\).
The three distinct lattices \(\Lambda_i\) are Leech lattices, every
pair has intersection \(J\), and
\[
 rJ=J,\qquad [\Lambda_i:J]=9,\qquad\det J=81.
\]
Put \(u=(\omega_2,\omega_4,0,\omega_2;0)\),
\(v=\rho/6-\alpha\), and \((e_1,e_2,e_3,e_4)=(u,v,ru,rv)\).
Their classes form a basis of \(J^\vee/J\) over \(\mathbb F_3\).
In this basis,
\begin{equation}\label{eq:B-disc}
 \bigl((e_i,e_j)\bigr)=
 \begin{pmatrix}
 4&2&-2&-4/3\\2&4&-2/3&-2/3\\
 -2&-2/3&4&2\\-4/3&-2/3&2&4
 \end{pmatrix},\qquad
 Q_J(a,b,c,d)=\frac{2ad+bc+bd}{3}\pmod{\Z}.
\end{equation}
In particular this discriminant form is split and nondegenerate.
\end{statement}

\begin{proof}
We give both the exact finite verification and the algebra that makes
it sufficient. The root lattice has determinant \(6^4\cdot4=5184\).
The classes in \eqref{eq:B-glue} span an isotropic subgroup of its
quadratic discriminant group of order \(72\): reduce the four
\(A_5\) coordinates modulo six and the \(D_4\) classes modulo its
root lattice. The orders of the independent primary components are
\(2^3\) and \(3^2\). Their pairings and half norms are integral.
Consequently \(N\) is even and has determinant \(5184/72^2=1\).

Here and below an assertion about a lattice span is verified by integer
Hermite normal form in the ambient coordinates multiplied by twelve.
The complete implementation is \path{code/B_exact_lattices.py};
the constructed bases and its output are in the correspondingly named
\path{records/B_*} files. A basis of \(K_d\) is obtained by taking
the Hermite normal form of
\[
 d\alpha,\qquad g-(\rho,g)\alpha
 \quad(g\text{ a generator of }N).
\]
The homomorphism \(x\mapsto(\rho,x)\pmod d\) is onto since it takes
\(\alpha\) to one, so \([N:K_d]=d\). The added vector \(w_d\)
in \eqref{eq:B-neighbor} pairs integrally with \(K_d\), has square
\(4(6/d)^2\), and has exact order \(d\) modulo \(K_d\).
These facts prove evenness and unimodularity of \(N_d\).

For clarity, intersection is not inferred from the coincidence of any
shells. The identity
\[
 A\cap B=(A^\vee+B^\vee)^\vee
\]
for full lattices is applied with exact rational Gram inverses and
Hermite normal forms. It gives determinant \(81\) for all three pairwise
intersections, identical bases for those intersections, and invariance
under \(r\). The determinant formula for an inclusion in a unimodular
lattice gives index nine. Direct reduction gives \(3e_i\in J\),
\((e_i,J)\subset\Z\), and
\(J+\sum_i\Z e_i=J^\vee\). The last lattice has index \(81\), so the
four classes are independent. Taking half the norm in the displayed
Gram matrix gives \eqref{eq:B-disc}. The plane spanned by its first two
basis vectors is totally isotropic. A nondegenerate four-dimensional
quadratic space containing such a plane is split.

We verify root systems by complete short-vector enumeration, not by a
random or a numerical search. The exact enumeration method is described
below. It gives \(144\) roots for \(N\), the roots already in \(R\),
and no roots in any \(\Lambda_i\). The uniqueness of the rootless even
unimodular positive definite lattice of rank 24 then identifies these
as Leech lattices. This is the rank-24 classification, in the form
stated and proved in Chapter II, Corollary 3.14, of
\citet{B_ChenevierLannes2019}; the rootless case is also established
by the construction and uniqueness argument of \citet{B_Borcherds1985}.
Every hypothesis here has been checked: rank 24, positive definiteness,
evenness, determinant one, and the complete absence of roots.
\end{proof}

\begin{statement}[B.2 (Proved): all eight even unimodular extensions]
There are exactly eight even unimodular lattices containing \(J\) in
\(J\otimes\Q\). In the coordinates of \eqref{eq:B-disc} they are
\(M_i=J+\Z\widetilde a_i+\Z\widetilde b_i\), where a tilde denotes the
linear combination of the \(e_j\) with the indicated coefficients:
\begin{center}\small
\begin{tabular}{ccllc}
\toprule
\(i\)&\(a_i\)&\(b_i\)&Root system&\(rM_i\)\\\midrule
0&0001&0010&none; \(\Lambda_1\)&\(M_7\)\\
1&0001&1100&none&\(M_3\)\\
2&0010&1000&\(A_2^{12}\); \(N_2\)&\(M_2\)\\
3&0011&1200&none&\(M_4\)\\
4&0012&0100&none&\(M_1\)\\
5&0100&1000&none; \(\Lambda_0\)&\(M_0\)\\
6&0112&1020&\(A_2^{12}\)&\(M_6\)\\
7&0121&1010&none; \(\Lambda_2\)&\(M_5\)\\\bottomrule
\end{tabular}
\end{center}
The two rooted extensions have exactly the same set of 72 roots in the
ambient Euclidean space. Moreover \(N_3\) has root system \(A_1^{24}\).
Thus Lead 6 holds, including its assertions about the action of \(r\).
\end{statement}

\begin{proof}
An intermediate lattice \(J\subset M\subset J^\vee\) is even exactly
when \(M/J\) is isotropic for \(Q_J\); it is unimodular exactly when
\((M/J)^\perp=M/J\). Indeed integrality and evenness are respectively
the vanishing of the induced bilinear form and of the quadratic form,
and the dual corresponds to the perpendicular subgroup. In our
nondegenerate group of order \(3^4\), these are precisely the isotropic
planes. The split four-space has 32 nonzero isotropic vectors and
16 isotropic lines. Every isotropic line belongs to two maximal
isotropic planes and each plane contains four lines, hence there are
\(16\cdot2/4=8\) planes. Alternatively substitution in
\(2ad+bc+bd=0\) gives exactly the eight spans in the table; the program
checks all 81 vectors and all pairs. The two rulings each have four
planes, as follows also by fixing one plane and using the parity of
the intersection dimension.

For each span the program computes an even unimodular basis and
exhausts every vector of square at most two. Six counts are zero;
the remaining counts are 72. Joining roots when their inner product
is nonzero partitions each of the latter root sets into twelve
rank-two components of size six, which are \(A_2\). Exact equality
of the two ambient root sets is checked. The same enumeration for
\(N_3\) gives 48 roots in 24 orthogonal rank-one components.
Applying \(r\) to the bases and reducing their spans gives the last
column. This verifies both fixed rooted lattices and the two cycles
of length three on the Leech lattices.
\end{proof}

\paragraph{The exact enumeration certificate.}
The program \path{code/B_enumerate.cpp} uses an integral,
fraction-free version of completing squares. After an integral
unimodular basis reduction, whose change-of-basis matrix and determinant
are checked, let \(G\) be the Gram matrix and \(D_i\) its first \(i\)
principal minors, with \(D_0=1\). Rational elimination gives integers
\(m_{ij}\) such that
\[
 x^tGx=\sum_{i=0}^{23}
 \frac{(D_{i+1}x_i+\sum_{j>i}m_{ij}x_j)^2}{D_iD_{i+1}}.
\]
All \(D_i\) are positive. Enumerate backwards from \(i=23\), using
integer square roots for the bounds. If the residual after the later
coordinates is \(R_{i+1}/D_{i+1}\), the new integral residual is
\[
 R_i=\frac{D_iR_{i+1}-(D_{i+1}x_i+\sum_{j>i}m_{ij}x_j)^2}{D_{i+1}}.
\]
The division is exact by the completed-square identity. Every permitted
integer in the closed interval is visited. This gives an induction
proof of completeness, including points on its boundary. The files
\texttt{*.enum} contain the exact minors and coefficients; files
\texttt{*.basis.json} identify the reduced bases. Root enumerations
also record every vector. No floating point decision enters the
enumeration. The largest minor in the construction of Statements
B.1--B.2 is \(2596\).

\subsection{Minimum, theta series, and the full automorphism group}

\begin{statement}[B.3 (Proved): the theta series of \(J\)]
The minimum of \(J\) is four, its kissing number is \(23058\), and
its theta series is a modular form of weight 12 on \(\Gamma_0(3)\)
with trivial character. With the usual normalizations of \(E_4\),
\(\Delta=\eta^{24}\), and \(\eta\), it is
\begin{equation}\label{eq:B-theta}
\begin{split}
 \Theta_J(\tau)={}&\frac{7381}{66430}E_4(\tau)^3
 +\frac{59049}{66430}E_4(3\tau)^3
 -\frac{531432}{6643}\Delta(\tau)\\
 &-\frac{4251528}{6643}\Delta(3\tau)
 +\frac{88938}{73}\eta(\tau)^{12}\eta(3\tau)^{12}.
\end{split}
\end{equation}
In particular its coefficients through \(q^8\) are
\begin{align*}
 1+23058q^2+1849032q^3+44291124q^4
 &+514246104q^5+3824355978q^6\\
 &+20831769864q^7+90541677312q^8+O(q^9).
\end{align*}
\end{statement}

\begin{proof}
The exact enumeration gives respectively
\(1,0,23058,1849032,44291124\) vectors of squares \(0,2,4,6,8\).
The norm-eight enumeration visits \(215418506\) nodes. Its complete
output and input are \path{records/B_J8.out} and
\path{records/B_J8.enum}; the separate norm-six run is retained.
Positive definiteness and the absence of norm-two vectors then prove
the first two assertions.

For the modular assertion we include the transformation argument.
Evenness gives \(\Theta_J(\tau+1)=\Theta_J(\tau)\). Poisson summation
on the rank-24 lattice gives
\[
 \Theta_J(-1/\tau)=\frac{\tau^{12}}9\Theta_{J^\vee}(\tau).
\]
Equation \eqref{eq:B-disc} gives \(3(x,x)/2\in\Z\) for every
\(x\in J^\vee\), so \(\Theta_{J^\vee}\) is invariant under
\(\tau\mapsto\tau+3\). The weight-12 slash action therefore makes
\(\Theta_J\) invariant under \(T\) and \(ST^3S^{-1}\), and hence
under \(\Gamma_0(3)\). To check the last generation assertion, write
the first column of a matrix as \((a,3c)^t\). Left multiplication by
these two parabolic generators permits \(a\mapsto a+3nc\) and
\(c\mapsto c+ma\). Nearest-integer reduction gives successively
\(|a|\le3|c|/2\) and \(|c|\le|a|/2\), reducing a nonzero
\(|c|\) by a factor at most \(3/4\). The terminal triangular matrix
is generated by \(T\) and \(-I\), the latter acting trivially in
weight 12. Holomorphy at the two cusps follows from the nonnegative
norms in the two Poisson expansions. The lattice level is exactly three, because \eqref{eq:B-disc} has
a value \(1/3\). The theta form itself is not a form for the full
modular group: its constant and linear coefficients would force
it to equal \(E_4^3-720\Delta\), whose quadratic coefficient is
196560 rather than 23058.

Each term on the right of \eqref{eq:B-theta} is holomorphic of weight
12 on \(\Gamma_0(3)\). For the eta product this follows directly from
\(\eta(\tau+1)=e^{\pi i/12}\eta(\tau)\) and
\(\eta(-1/\tau)=\sqrt{-i\tau}\eta(\tau)\): the product is invariant
under \(T\), and its Fricke transformation is
\(g(-1/(3\tau))=3^6\tau^{12}g(\tau)\), which gives the other
generator. It vanishes at both cusps.

Finally, a nonzero weight-12 holomorphic form on \(\Gamma_0(3)\)
cannot vanish to order more than four at infinity. One proof takes
the product of its four slash translates over
\(\Gamma_0(3)\backslash\mathrm{SL}_2(\Z)\); the product is a nonzero
holomorphic form of weight 48 for the full modular group. Its order
at infinity is at least that of the original form, and the usual
valence formula bounds that order by \(48/12=4\). The right side of
\eqref{eq:B-theta} has the five enumerated coefficients, so subtraction
and this bound prove the identity. The coefficient computation, using
integer products for eta and divisor sums for \(E_4\), is
\path{code/B_theta.py}; its five-by-five coefficient matrix has
nonzero determinant \(-387419760\). No extrapolation of shell counts
is used in the identity.
\end{proof}

The equality of the two root sets in Statement B.2 makes the group
calculation substantially smaller. Let \(R'\cong A_2^{12}\) be their
common span. In each component choose roots \(\beta_i,\gamma_i\)
with \((\beta_i,\gamma_i)=-1\), both having the same nonzero class
under the quotient homomorphism \(R'\to R'/(J\cap R')\cong\mathbb F_3\).
These choices, in ambient coordinates, are recorded in
\path{records/B_autJ_data.json}.

\begin{statement}[B.4 (Proved): the full group of \(J\)]
The full isometry group of \(J\) has order
\[
 |\operatorname{Aut}J|=12597120=2\cdot3^7\cdot2880
                      =2^7 3^9 5.
\]
An explicit structural description is
\[
 \operatorname{Aut}J\cong C_2\times
    \bigl((\mathbb F_3)^7\rtimes\operatorname{Aut}_{\rm mon}(D)\bigr),
\]
where \(D\) is the ternary code generated by
\begin{equation}\label{eq:B-code}
\begin{pmatrix}
1&1&0&0&0&0&0&0&2&2&2&2\\
0&0&1&0&0&2&0&2&2&2&0&2\\
0&0&0&1&0&2&0&0&2&1&2&1\\
0&0&0&0&1&2&0&2&1&0&1&1\\
0&0&0&0&0&0&1&2&1&1&2&2
\end{pmatrix}.
\end{equation}
Here \(\operatorname{Aut}_{\rm mon}\) means signed coordinate
permutations. Its order is 2880. Fourteen integral matrices generating
the full group, in the basis \path{B_J_basis.json}, are supplied in
\path{B_autJ_data.json} and verified by \path{B_autJ.py}.
Its action on the eight extensions has order 36 and has two orbits,
the six Leech extensions and the two rooted extensions.
\end{statement}

\begin{proof}
Every isometry of \(J\) permutes its even unimodular extensions and
preserves roots. It therefore preserves \(R'\), intrinsically obtained
from the two rooted extensions. This gives an upper bound by root
system isometries, and is the reason that the computation below
identifies the full group.

Put \(K=J\cap R'\). Exact reduction of the displayed root basis gives
\[
 K=\{\textstyle\sum_i(a_i\beta_i+b_i\gamma_i):
       a_i,b_i\in\Z,\ \sum_i(a_i+b_i)=0\pmod3\}.
\]
For \(x=\sum_i(a_i\beta_i+b_i\gamma_i)\in R'^\vee\), put
\(c_i=3b_i\pmod3\). Then \(a_i+b_i\in\Z\), and the homomorphism
\[
 R'^\vee/K\longrightarrow\mathbb F_3^{12}\oplus\mathbb F_3,
 \quad x\longmapsto
 \left(c,\sum_i(a_i+b_i)-\sum_i c_i\right)
\]
is an isomorphism; representatives of \(c_i\) can be changed by
multiples of three without changing the last coordinate. Reduction of
the 24 generators of \(J\) gives the graph of a functional \(f:D\to
\mathbb F_3\), with values \((0,1,1,2,1)\) on the rows of
\eqref{eq:B-code}. The code has weight enumerator
\(1+90z^6+140z^9+12z^{12}\).

In one component let \(r_i\) send \(\beta_i\) to \(\gamma_i\)
and \(\gamma_i\) to \(-\beta_i-\gamma_i\). Let \(t_i\) be the
involution with row-coordinate matrix
\(\left(\begin{smallmatrix}0&1\\1&0\end{smallmatrix}\right)
\left(\begin{smallmatrix}0&1\\-1&-1\end{smallmatrix}\right)\).
In the quotient coordinates above, \(r_i\) sends
\((c,h)\) to \((c,h-c_i)\), whereas \(t_i\) changes \(c_i\) to
\(-c_i\) and fixes \(h\). The product of the rotations with exponents
\[
 z_0=(0,0,1,1,2,0,1,0,0,0,0,0)
\]
changes the graph of \(f\) to \(D\oplus0\).

An isometry of \(A_2^{12}\) preserving \(K\) must act on its defining
nonzero linear functional by one common scalar \(\varepsilon=\pm1\).
The component root isometries are the dihedral groups of order twelve.
Those with scalar \(+1\) are generated by \(r_i,t_i\); the common
scalar \(-1\) is provided by the central global negation. After the
normalizing rotations, preservation of \(D\oplus0\) is therefore
equivalent to a monomial automorphism of \(D\), and a rotation vector
\(z\in D^\perp\). The latter space has dimension seven. This proves
the claimed group description, both its upper bound and existence.

We certify the finite code group twice. The first computation encodes
all 242 nonzero codewords and the 24 signed coordinates as a coloured
graph; the six returned signed coordinate generators are checked
against every codeword, and an exact permutation-group calculation
gives order 2880. A separate exhaustive computation in
\path{B_code_aut_exhaust.cpp} does not rely on the graph
algorithm. The five pivot columns of \eqref{eq:B-code} are the standard
basis of \(\mathbb F_3^5\). Its twelve columns form eleven projective
classes, one occurring twice. We exhaust all ordered images of the
five pivot columns among these eleven classes, with all \(2^5\) signs:
\(11\cdot10\cdot9\cdot8\cdot7\cdot32=1774080\) choices. Exactly
1440 linear transformations preserve the column multiset. Each has
two lifts to a signed coordinate permutation, arising from the repeated
class. The 2880 resulting permutations are distinct and all belong to
the group generated by the six displayed generators. This proves the
order without an unverified completeness assertion about a search.

The seven rotations from a basis of \(D^\perp\), the six monomial
generators, and global negation give fourteen generators. Conjugation
by the displayed normalizing rotations and by the exact root-to-\(J\)
basis matrix gives the recorded integral matrices. Each has determinant
\(\pm1\) and satisfies \(T G_JT^t=G_J\). Reducing their images of each
extension gives the stated permutation group and its two orbits.
\end{proof}

\begin{statement}[B.5 (Proved): the associated plane and its stabilizer]
Identify \(\Lambda=\Lambda_0=M_5\). The plane
\[
 P=3J^\vee/3\Lambda\ \subset\ \Lambda/3\Lambda
\]
is totally isotropic and has the two generators \(3ru,3rv\pmod{3\Lambda}\).
Its perpendicular preimage in \(\Lambda\) is \(J\). Its full
stabilizer in \(\mathrm{Co}_0=\operatorname{Aut}\Lambda\) has order
\(2099520\). Its \(\mathrm{Co}_0\)-orbit consists of
\(3960692736000\) isotropic planes. Intrinsically, it is the orbit
whose perpendicular lattice is isometric to the \(J\) above; the
choice of a Leech extension introduces no further orbit, because
\(\operatorname{Aut}J\) is transitive on all six such extensions.
Eleven integral generators in the displayed \(\Lambda_0\) basis are
given in \path{records/B_specific_plane_group.json}. This is orbit 4
in the complete classification below; the integral comparison
isometry is supplied in \path{B_specific_to_orbit4.json}.
\end{statement}

\begin{proof}
Since \(3J^\vee\subset J\subset\Lambda\), the quotient is defined.
The image has order nine: \([J^\vee:\Lambda]=9\). Pairings among
its vectors are divisible by three and its norms are divisible by
six, by \eqref{eq:B-disc}, proving isotropy for
\(x\mapsto x^2/2\pmod3\). The equality
\[
 x\in J\quad\Longleftrightarrow\quad
 (x,J^\vee)\subset\Z
 \quad\Longleftrightarrow\quad
 (x,3J^\vee)\subset3\Z
\]
proves the perpendicular assertion. The first two discriminant
basis classes generate \(\Lambda/J\), so the last two give the
specified basis of \(P\).

The stabilizer equals \(\operatorname{Aut}(J)_\Lambda\): in one
direction an isometry preserving \(P\) preserves its perpendicular
preimage; in the other it preserves \(3J^\vee/3\Lambda\). Statement
B.4 and orbit--stabilizer give \(12597120/6=2099520\). The order
\(8315553613086720000\) of \(\mathrm{Co}_0\), twice the listed
order of \(\mathrm{Co}_1\) in the primary ATLAS table
\citep{B_ATLASCo1}, then gives the displayed
orbit length. The isometry class characterization follows because
an isometry between two copies of \(J\) can be adjusted by an
automorphism of \(J\) to match any prescribed Leech extensions.
For the generator assertion we apply Schreier's lemma to the
six-point orbit of \(M_5\) under the fourteen matrices of B.4.
If \(t_i\) takes \(M_5\) to \(M_i\), the words
\(t_i g t_{g(i)}^{-1}\) generate its stabilizer. After change of
basis to \(\Lambda_0\), every retained matrix is integral,
unimodular, and preserves its Gram matrix. Their faithful action on
the 72 intrinsic roots has independently computed order 2099520;
\path{code/B_specific_plane_group.py} retains the words and checks.
\end{proof}

\subsubsection{Complete classification of isotropic planes modulo three}

For an isotropic line \(\ell\subset\Lambda/3\Lambda\), let
\(N(\ell)\) denote its unique even unimodular 3-neighbour. Explicitly,
choose \(v\in\Lambda\) spanning \(\ell\), replace \(v\) by
\(v+3a\) so that \(v^2\equiv0\pmod {18}\), and put
\[
 N(\ell)=\{x\in\Lambda:(x,v)\equiv0\pmod3\}+\Z(v/3).
\]
The replacement exists because the nonzero functional
\((v,{\cdot})\pmod3\) is onto. The resulting lattice is independent
of the replacement and sign of \(v\). Indeed the common index-three
intersection has split discriminant plane; its two isotropic lines
give precisely \(\Lambda\) and \(N(\ell)\). We call \((\ell,P)\),
where \(\ell\subset P\) is a totally isotropic plane, a marked plane.

\begin{statement}[B.9 (Proved): all 129 plane orbits and their full groups]
There are exactly 129 \(\mathrm{Co}_0\)-orbits of totally isotropic
planes in \(\Lambda/3\Lambda\). Of these, 106 have a line whose
3-neighbour is \(A_2^{12}\), another 20 have a line whose
3-neighbour is \(A_1^{24}\), and the remaining three have eight
Leech extensions. Table~\ref{tab:B129} gives every orbit, its full
stabilizer order, its orbit length, and the eight Niemeier extensions
of the perpendicular preimage \(J_P\).

For every row we supply a two-by-24 reduced matrix over \(\mathbb F_3\)
in the same integral Leech basis, not merely an abstract isometry
class. These matrices and all eight extension labels are in
\path{records/B_all_129_plane_orbits.json} and its CSV companion.
Full integral stabilizer generators in that same basis, with exact
order certificates, are in \path{records/B_plane_groups/001.json}
through \path{records/B_plane_groups/129.json}. The sum of the 129
orbit lengths is exactly
\[
 61549107892647466000,
 \qquad
 \sum_{[P]}\frac1{|\mathrm{Co}_{0,P}|}
   =\frac{14703561369481}{1986515435520}.
\]
The eight-extension multiset alone does not determine the orbit:
rows 1 and 12 both have four Leech and four \(A_2^{12}\) extensions,
but their stabilizers have orders 93312 and 103680.
\end{statement}

\begin{proof}
We give the finite reductions, their completeness arguments, and the
independent mass check. All quadratic forms in the finite calculation
are \(x\mapsto x^2/2\pmod3\).

The binary Golay code used in the common coordinate model is the
span of the following twelve hexadecimal masks; bit \(i\) means
coordinate \(i\), starting at zero:
\[
\begin{gathered}
\mathtt{fff0f0},\ \mathtt{ff0ff0},\ \mathtt{f0fff0},\
\mathtt{0ffff0},\ \mathtt{cccc00},\ \mathtt{aaaa00},\\
\mathtt{6ac0c0},\ \mathtt{c6a0a0},\ \mathtt{a6c00c},\
\mathtt{6ca00a},\ \mathtt{11111e},\ \mathtt{ffffff}.
\end{gathered}
\]
In \(\R^{24}\) with scalar product divided by eight, take the
integer span of \(8e_0\), the 23 vectors \(4(e_i-e_0)\), twice
these twelve binary code generators, and \((-3,1^{23})\).
Exact Hermite reduction gives the common basis; its Gram matrix is
even of determinant one, and exact norm-two enumeration is empty.
This identifies it as the Leech lattice. The ternary glue code in
the rooted marked model is the row span over \(\mathbb F_3\) of
\[
\begin{pmatrix}
1&0&0&0&0&1&0&2&2&1&1&0\\
0&1&0&0&0&2&0&1&0&1&1&2\\
0&0&1&0&0&2&0&2&2&2&0&2\\
0&0&0&1&0&2&0&0&2&1&2&1\\
0&0&0&0&1&2&0&2&1&0&1&1\\
0&0&0&0&0&0&1&2&1&1&2&2
\end{pmatrix}.
\]
Its self-duality and weight distribution are checked directly on
all 729 words. These displayed codes specify the inputs independently
of any software convention for Golay coordinates.

\paragraph{The four line orbits and full line stabilizers.}
There are four line types. The familiar labels \(3,6_{22},6_{23},9_{33}\)
refer to vector classes, not to their dimension. We use them only as
names for the explicit representatives in the supplied programs.
Their full line-stabilizer orders and orbit lengths are
\[
\begin{array}{c|r|r|c}
\text{type}&|\mathrm{Co}_{0,\ell}|&|\mathrm{Co}_0\ell|&N(\ell)\\\hline
3&991533312000&8386560&\Lambda\\
6_{22}&110361968640&75348000&\Lambda\\
6_{23}&489646080&16982784000&A_1^{24}\\
9_{33}&277136640&30005248000&A_2^{12}
\end{array}
\]
We justify fullness rather than inferring it from a generated subgroup.
For the last two rows, work in the indicated rooted neighbour \(N\).
The line stabilizer is exactly the group of root-system isometries
preserving both its glue and \(K=N\cap\Lambda\). Preservation of
\(N,K\) preserves the other unimodular extension \(\Lambda\), by
the two-line discriminant argument above.

For \(A_1^{24}\), in orthogonal root coordinates,
\(N=\Z^{24}+\tfrac12 C_{24}\), with the extended binary Golay
code, and \(K=\ker(\rho,{\cdot})\pmod3\),
\(\rho=\tfrac12\sum_i\alpha_i\). A signed permutation preserving
\(K\) must have a common sign on all roots. Thus the full group is
\(C_2\times M_{24}\), of order 489646080. The code graph and the
displayed coordinate permutations independently give the full
\(M_{24}\) order 244823040.
For \(A_2^{12}\), its glue is the self-dual ternary Golay code
\(C_{12}\) of dimension six. In a root basis
\(\beta_i,\gamma_i\) with \(\rho=\sum_i(\beta_i+\gamma_i)\),
the same quotient-functional argument used in B.4 gives
\[
 \mathrm{Co}_{0,\ell}
   =C_2\times\bigl(3^6\rtimes
           \operatorname{Aut}_{\rm mon}(C_{12})\bigr).
\]
The rotations are indexed by \(C_{12}^{\perp}=C_{12}\). The
monomial group, computed on the complete code graph, has order
190080, giving the displayed full order 277136640. All its lifted
24-dimensional matrices, including the stabilizers of the four
codeword strata below, are retained.

For the first two rows we also provide independent finite upper
bounds. Use the standard integral-coordinate Leech model, with
scalar product \(x\cdot y/8\), in
\path{B_standard_Leech_basis.json}. For type 3 take
\(v=(5,1^{23})\), of squared norm six. It is the unique vector of
norm six in its residue class: two such vectors differing by
\(3\Lambda\) would have squared distance at least 36 and at most
24. The 552 norm-four vectors \(w\) with \((v,w)=3\) span the
ambient space. Their Gram values off the diagonal are \(-1,1,2\);
the \(-1\) pairs are \(\{w,v-w\}\). We encode the inner-product-one
graph and separately mark those 276 pairs. Its full permutation
group has order 991533312000. Exact reconstruction of its linear
maps exhibits a generator that fails to preserve the Leech lattice.
The independently verified lattice subgroup of order 495766656000
has index two in that graph group. The lattice vector stabilizer
therefore has precisely that order, and adjoining global negation
gives the stated line stabilizer. The graph, maps, and failed
integrality witness are in \path{B_Co3_upper_bound.json}.

For type \(6_{22}\), take the three vectors whose nonzero coordinate
triples are the permutations of \((-8,4,4)\), in coordinate
positions 2,3,4. They have squared norm 12, sum zero, and lie in the
same nonzero class modulo \(3\Lambda\). These are all the shortest
representatives of that class. Distinct representatives of norm 12
must have inner product \(-6\): their difference has squared norm
at least 36, and the next possible squared norm in \(3\Lambda\)
is 54, larger than the maximum 48. A set of \(k\) such vectors has
Gram matrix \(18I_k-6\mathbf1\mathbf1^t\), whence \(k\le3\).
Their differences divided by three form a triangle of minimal
vectors. Fix two of these minimal vectors \(a,b\), with norms four
and inner product \(-2\). The 891 minimal vectors satisfying
\((w,a)=2,(w,b)=0\) span \(b^\perp\). Their Gram graph has full
group order 55180984320. A verified pointwise triangle subgroup has
order 9196830720. Of the six cosets of that subgroup in the graph
group, exact extension fixing \(b\) shows that just one preserves
the Leech lattice. Consequently the full pointwise triangle group
has order 9196830720. The displayed generators induce all six
permutations of the three norm-12 representatives, and global
negation doubles the class stabilizer to the line stabilizer.
This proves the second row. The complete six-coset test is
\path{B_U6_upper_bound.py}; no untested graph symmetry is counted as
a lattice isometry.

These calculations independently recover the customary names
\(\mathrm{Co}_3\) and \(U_6(2):S_3\). The corresponding known
stabilizer interpretations are recorded in
\citet[Section 2]{B_BrouwerLeech} and
\citet[Section 5, p.~10, statement preceding Corollary 7]{B_Fairbairn2010}.
Fairbairn's hypothesis is an unordered triangle of three minimal
vectors with sum zero; the difference triangle just constructed
has exactly that form. We use our graph and lattice calculations
for the numerical upper bounds. The first two types are distinct
because one has a norm-six representative and the other does not;
the latter assertion also follows from the minimum-distance bound
between a norm-six and a norm-12 vector in one class. The last two
types are distinguished by their rooted neighbours. Finally
\[
 8386560+75348000+16982784000+30005248000
 =47071766560
 =\frac{(3^{11}+1)(3^{12}-1)}2.
\]
This is the number of isotropic lines in the split quadratic
24-space. Thus the four verified line orbits exhaust all lines.

\paragraph{All planes with an \(A_2^{12}\) neighbour.}
For this line the quotient \(\ell^\perp/\ell\) is explicitly
\[
 W=K/(3N+\Z\rho),\qquad\dim_{\mathbb F_3}W=22.
\]
The comparison sends a vector of \(K\) to the same rational vector
in the Leech construction. Its kernel is \(3N+\Z\rho\); the
displayed matrices check the pairing, so it is an isometry of
quadratic spaces. Projection to the glue gives an equivariant
map \(W\to C_{12}\) with 16-dimensional kernel. The four codeword
orbits have weights \(0,6,9,12\) and sizes \(1,264,440,24\).
For a fixed representative of each, the full stabilizer acts
affinely on its 16-dimensional fibre. We give all affine matrices,
their integral lattice lifts, and the quadratic equation in
\path{B_A2_strata_setup.json}. The four complete enumerations give
\[
\begin{array}{c|r|r|r}
\text{codeword weight}&\text{code orbit}&
 \text{isotropic fibre vectors}&\text{fibre orbits}\\\hline
0&1&14467005&16\\
6&264&14348907&54\\
9&440&14348907&64\\
12&24&14348907&14
\end{array}
\]
Each enumeration examines all \(3^{16}=43046721\) vectors and
closes every retained orbit under every generator. Removing the
zero vector gives 147 marked-plane orbits. Since global negation
belongs to the full line group, a fibre orbit of size \(f\) over
a codeword orbit of size \(c\) gives a marked-plane orbit of size
\(fc/2\). Their sum is exactly
\[
 5230235650
   =\frac{(3^{10}+1)(3^{11}-1)}2,
\]
the full number of isotropic lines in \(W\). The dense orbit maps
are supplied losslessly compressed, together with their hashes.

For each of the 147 representatives we construct \(J_P\), its
split four-dimensional discriminant form, and all eight maximal
isotropic subspaces, hence all eight unimodular extensions. Complete
root enumeration identifies each extension. The four extensions
meeting the original Leech lattice in index three correspond
bijectively to the four lines of \(P\); the remaining three other
extensions have intersection index nine. This distinction is
checked by exact lattice intersections, not inferred from a root
label.

We change the marked line in every possible way that again gives
an \(A_2^{12}\) neighbour. In its root system choose
\(\beta_i,\gamma_i\) on which the intersection functional has
value one. Root-coordinate rotations normalize its value on the
glue, and an isomorphism of the complete Golay code graphs maps
the glue to the fixed model. The resulting rational map \(T\)
satisfies both the exact Gram identity and
\(\Lambda T=\Lambda\). Reducing the transported second vector in
the fibre tables gives the new marked orbit. These are actual
isometries, rather than matches of extension multisets. Their
connected components are exactly the unmarked plane orbits:
any isometry between two planes first sends the chosen mark to one
of the four possible marks, after which it lies in the full line
stabilizer already enumerated. The 147 marked orbits fuse to 106
unmarked orbits. All links and representatives are in
\path{B_A2_plane_orbits.json}.

\paragraph{All additional planes with an \(A_1^{24}\) neighbour.}
In the orthogonal root coordinates the corresponding quotient is
\[
 W=\{x\in\mathbb F_3^{24}:\sum_i x_i=0\}/\langle\mathbf1\rangle.
\]
An isotropic vector has all three colour-class sizes divisible by
three. Quotienting by constants and nonzero scalars makes these
three classes unordered. Apart from the zero vector, their sorted
sizes are
\[
\begin{gathered}
(0,3,21),(0,6,18),(0,9,15),(0,12,12),\\
(3,3,18),(3,6,15),(3,9,12),(6,6,12),(6,9,9).
\end{gathered}
\]
We first fix the smallest class. Its size is 0, 3 or 6; there are
four \(M_{24}\)-orbits, with lengths
\(1,2024,21252,113344\) and stabilizers of orders
\(244823040,120960,11520,2160\). The last two lengths sum to
\(\binom{24}{6}\). Exhaustive subset orbits of the second class,
followed by all allowed permutations of the three colours, give
114 marked-plane orbits. Their lengths again sum to 5230235650.
Every extension of every representative is constructed and its
roots are enumerated. Of these marked orbits, 83 have an
\(A_2^{12}\) mark and are mapped by the preceding explicit
isometry construction into the 106 existing classes. The other
31 marked orbits fuse, by all \(A_1^{24}\) changes of mark, to
20 further classes. The binary Golay graph gives the exact
coordinate permutation in every such change. The data and all
cross-identifications are in \path{B_A1_plane_orbits.json}.

\paragraph{The last three classes.}
The 126 classes just obtained contain
\(61549106976378112000\) planes. The residual count is
\(916269354000\). We find three further representatives in the
standard Leech coordinates; their second vectors are recorded by
the following unambiguous packing convention:
\[
 \operatorname{pack}(w)=
   \sum_{w_i=1}2^i+2^{24}\sum_{w_i=2}2^i,
 \quad w_i\in\mathbb F_3,\quad 0\le i<24.
\]
For type 3 the anchor is \((5,1^{23})\), and for type \(6_{22}\)
it is the vector with entries \((-8,4,4)\) in positions 2,3,4.
The table gives the complete orbit under the anchor-class group
(global negation acts trivially on marked planes).
\[
\begin{array}{c|r|r|r|r}
\text{anchor}&\operatorname{pack}(w)&\text{flag orbit}&
 |\mathrm{Co}_{0,P}|&|\mathrm{Co}_0P|\\\hline
3&9192064&37950&52254720&159134976000\\
6_{22}&416910284928&27720&15925248&522161640000\\
6_{22}&5174966569752&12474&35389440&234972738000
\end{array}
\]
The first plane has two type-3 lines and two type-\(6_{22}\) lines;
the other two have four type-\(6_{22}\) lines. Every eligible change
of mark stays in the same displayed flag orbit. All changes are
checked as integral Leech isometries, and their complete orbit
trees, parent edges and generators are retained. Thus these are
three distinct full plane orbits. Exact enumeration shows that
all 24 of their extensions are rootless. The seeded search used
to discover these representatives is not a completeness argument:
completeness follows from their closed group orbits and the exact
mass equality below.

\paragraph{Full group generators and order certificates.}
For clarity we state the constructive recipe used for every row.
In a rooted marked model let \(H=\mathrm{Co}_{0,\ell}\), and let
\(H_P\) be the stabilizer of the marked plane. In the \(A_2\)
model write \(H=C_2\times(R\rtimes M)\), where
\(R=C_{12}\cong3^6\) is the root-rotation group and
\(M=\operatorname{Aut}_{\rm mon}(C_{12})\). For a second quotient
vector \(y\), enumerate its at most 729 images under \(R\),
retaining a rotation \(r_x\) with \(yr_x=x\), and the rotation
kernel \(R_y\). For each signed monomial element \(m\) whose
image \(ym\) lies in that rotation orbit, retain
\(m r_{ym}^{-1}\). These elements and \(R_y\) generate the full
vector stabilizer by the semidirect decomposition. Adjoin global
negation to obtain \(H_P\). This gives a finite exact recipe using
only 729 rotations and 380160 signed monomial elements; it does
not enumerate a huge \(\mathrm{Co}_0\)-orbit. The program reduces
the generating set in its faithful action on the 72 roots and
checks that its order equals the independently known
\(|H|/|H\cdot P|\).

In the \(A_1\) model, adjoin to the 24 coordinate vertices and
759 octad vertices three vertices for the unordered colour classes.
The full coloured-graph group is exactly the \(M_{24}\) partition
stabilizer: each returned permutation is checked against every
octad and the three classes. Adjoin global negation. For either
model, for each mark whose transported flag lies in the original
\(H\)-orbit, compose the mark-changing isometry with an
\(H\)-transporter back to the representative. Adjoining these
matrices gives the full plane stabilizer, since its orbit on the
chosen mark has then been exhausted. Every generator is checked
to be integral in the common Leech basis, to preserve its Gram
matrix, and to preserve the two-dimensional plane modulo three.
Finally the union of the roots in all eight extensions spans the
ambient space and is intrinsic to \(J_P\). Its permutation action
is faithful; the orders of all 126 generated groups are independently
recomputed on that union. The complete program is
\path{B_plane_stabilizer_recipe.py}.

For the three remaining rows we use the saved flag-orbit trees.
If \(t_x\) is the tree word taking the representative to \(x\),
the words \(t_x g t_{xg}^{-1}\) generate its class-flag stabilizer.
Its order is computed in the faithful degree-276 or degree-891
class-group representation, giving respectively
\(13063680,1990656,4423680\). Adjoining global negation doubles
these orders. The exact changes of mark have anchor orbits of
sizes \(2,4,4\), giving precisely the three full orders displayed
above. All tree edges are rechecked by
\path{B_rare_plane_groups.py}; this proves both generation and
fullness without a guessed abstract group name.

\paragraph{Common coordinates and the final mass equality.}
Starting from the standard integral Leech basis, we construct its
type-\(9_{33}\) and type-\(6_{23}\) neighbours and use the root/glue
isometries above to identify them with the two marked models. The
exact maps preserve the forms and send the whole Leech lattice
onto the marked Leech lattice, as checked by integral Hermite
normal forms. Their inverses put every representative and every
group generator in one common basis. These comparison maps are
in \path{B_global_model_isometries.json}.

Independently of all orbit enumerations, an isotropic line has a
22-dimensional split perpendicular quotient, and a plane contains
four lines. Hence the total number of planes is
\[
 \frac14\,
 \frac{(3^{11}+1)(3^{12}-1)}2\,
 \frac{(3^{10}+1)(3^{11}-1)}2
 =61549107892647466000.
\]
The 126 rooted classes and the three new classes are disjoint by
the exhaustive changes of mark and their line types. Their exact
orbit lengths sum to this number: the last three lengths sum to
916269354000, precisely the residual count. Thus no plane remains.
Equivalently, division by
\(|\mathrm{Co}_0|=8315553613086720000\), as in
\citet[Theorem 1.1]{B_BrouwerLeech}, gives the stated rational mass.
The final stdout certificate \path{B_global_assembly.out} prints
both \texttt{RESIDUAL PLANES: 0} and \texttt{RESIDUAL MASS: 0}.
As a separate final check, \path{B_final_group_audit.py} rereads the
129 group files, verifies all 654 integral generator matrices, and
recomputes the same rational mass from their separately generated
full-group orders. Its stdout again gives both residuals as zero;
the per-file hashes are retained in \path{B_plane_group_manifest.json}.
\end{proof}

In Table~\ref{tab:B129}, the extension symbols have the following
fixed meanings; an exponent is the multiplicity of that entire
Niemeier lattice among the eight extensions, not a root multiplicity:
\[
\begin{array}{c|cccccc}
\text{symbol}&L&1&2&3&4&5\\\hline
\text{lattice}&\Lambda&A_1^{24}&A_2^{12}&A_3^8&A_4^6&A_5^4D_4
\end{array}
\qquad
\begin{array}{c|ccccc}
\text{symbol}&6&7&8&9&10\\\hline
\text{lattice}&D_4^6&A_6^4&A_7^2D_5^2&A_8^3&E_6^4
\end{array}
\]
The stratum and flag identifiers in the representative column refer
to the complete affine or subset orbit files described above. The
last three rows use the packed-coordinate convention in the proof.
The JSON and CSV also list all eight extensions separately, with
their exact bases and complete root sets referenced by filename.
{\scriptsize\setlength{\tabcolsep}{3pt}
\begin{longtable}{rllrrl}
\caption{All 129 isotropic-plane orbits. The representative is the stratum and flag identifier, or the packed rare vector; all planes are also supplied in one common integral Leech basis.}\label{tab:B129}\\
ID&mark&representative&$|\mathrm{Co}_{0,P}|$&$|\mathrm{Co}_0P|$&extensions\\\hline\endfirsthead
ID&mark&representative&$|\mathrm{Co}_{0,P}|$&$|\mathrm{Co}_0P|$&extensions\\\hline\endhead
1&A2&0,2&93312&89115586560000&$L^{4},2^{4}$\\
2&A2&0,3&23094720&360062976000&$L^{6},2^{2}$\\
3&A2&0,4&2592&3208161116160000&$L^{2},1^{3},2^{2},5$\\
4&A2&0,5&2099520&3960692736000&$L^{6},2^{2}$\\
5&A2&0,6&23328&356462346240000&$L^{2},2^{5},9$\\
6&A2&0,7&5038848&1650288640000&$L^{3},10,2^{4}$\\
7&A2&0,8&1920&4331017506816000&$L^{2},1^{2},2^{2},4^{2}$\\
8&A2&0,9&48&173240700272640000&$L,1,2^{5},4$\\
9&A2&0,10&12960&641632223232000&$L^{2},1^{3},2^{2},6$\\
10&A2&0,11&1296&6416322232320000&$L,1^{3},2^{4}$\\
11&A2&0,12&7776&1069387038720000&$L^{2},2^{6}$\\
12&A2&0,13&103680&80204027904000&$L^{4},2^{4}$\\
13&A2&0,14&648&12832644464640000&$L,2^{6},6$\\
14&A2&0,15&5184&1604080558080000&$L,2^{6},6$\\
15&A2&0,16&51840&160408055808000&$L^{2},2^{4},6^{2}$\\
16&A2&1,1&48&173240700272640000&$L,2^{6},5$\\
17&A2&1,2&4&2078888403271680000&$L,1,2^{5},4$\\
18&A2&1,3&12&692962801090560000&$L,2^{6},5$\\
19&A2&1,4&48&173240700272640000&$L,2^{5},3,4$\\
20&A2&1,5&4&2078888403271680000&$L,1,2^{4},3^{2}$\\
21&A2&1,6&4&2078888403271680000&$L,2^{5},3,4$\\
22&A2&1,7&4&2078888403271680000&$L,1^{2},2^{4},3$\\
23&A2&1,8&6&1385925602181120000&$L,2^{5},3,4$\\
24&A2&1,9&2&4157776806543360000&$L,1,2^{4},3^{2}$\\
25&A2&1,10&8&1039444201635840000&$L,1,2^{5},4$\\
26&A2&1,11&24&346481400545280000&$L,1,2^{5},4$\\
27&A2&1,12&4&2078888403271680000&$L,1,2^{4},3^{2}$\\
28&A2&1,13&4&2078888403271680000&$L,1^{2},2^{4},3$\\
29&A2&1,14&8&1039444201635840000&$L,1^{2},2^{4},3$\\
30&A2&1,15&36&230987600363520000&$L,2^{6},5$\\
31&A2&1,16&16&519722100817920000&$L,1^{2},2^{4},3$\\
32&A2&1,18&8&1039444201635840000&$L,1^{2},2^{4},3$\\
33&A2&1,19&48&173240700272640000&$L,1^{3},2^{4}$\\
34&A2&1,21&288&28873450045440000&$L^{2},2^{6}$\\
35&A2&1,23&12&692962801090560000&$L,2^{4},3^{3}$\\
36&A2&1,24&8&1039444201635840000&$L,1,2^{4},3^{2}$\\
37&A2&1,25&24&346481400545280000&$L,2^{4},3^{3}$\\
38&A2&1,27&48&173240700272640000&$L,2^{4},3^{3}$\\
39&A2&1,28&48&173240700272640000&$L,2^{4},3^{3}$\\
40&A2&1,29&32&259861050408960000&$L,1^{2},2^{4},3$\\
41&A2&1,30&48&173240700272640000&$L,1^{3},2^{4}$\\
42&A2&1,31&288&28873450045440000&$L^{2},2^{6}$\\
43&A2&1,32&24&346481400545280000&$L,1,2^{4},3,5$\\
44&A2&1,33&12&692962801090560000&$L,1^{2},2^{4},5$\\
45&A2&1,34&4&2078888403271680000&$L,1^{2},2^{3},3,4$\\
46&A2&1,35&16&519722100817920000&$L,1^{3},2^{3},4$\\
47&A2&1,37&16&519722100817920000&$L,1^{3},2^{2},3^{2}$\\
48&A2&1,38&64&129930525204480000&$L,1,2^{4},4^{2}$\\
49&A2&1,39&24&346481400545280000&$L,1^{3},2^{2},3^{2}$\\
50&A2&1,40&24&346481400545280000&$L,1^{3},2^{2},3^{2}$\\
51&A2&1,41&8&1039444201635840000&$L,1^{3},2^{2},3^{2}$\\
52&A2&1,42&16&519722100817920000&$L,1^{4},2^{2},3$\\
53&A2&1,43&48&173240700272640000&$L,1^{4},2^{2},3$\\
54&A2&1,44&192&43310175068160000&$L^{2},1^{2},2^{4}$\\
55&A2&1,46&8&1039444201635840000&$L,1^{3},2^{3},4$\\
56&A2&1,47&48&173240700272640000&$L,1,2^{5},7$\\
57&A2&1,48&480&17324070027264000&$L^{2},1^{2},2^{4}$\\
58&A2&1,49&24&346481400545280000&$L,1^{4},2^{2},3$\\
59&A2&1,50&384&21655087534080000&$L,1^{2},2^{4},8$\\
60&A2&1,51&120&69296280109056000&$L,1^{5},2,4$\\
61&A2&1,52&576&14436725022720000&$L^{2},1^{4},2^{2}$\\
62&A2&1,53&144&57746900090880000&$L,1^{4},2^{2},5$\\
63&A2&1,54&1440&5774690009088000&$L^{2},1^{4},2^{2}$\\
64&A2&2,1&4&2078888403271680000&$L,1^{2},2^{4},3$\\
65&A2&2,2&16&519722100817920000&$L,1^{4},2^{2},3$\\
66&A2&2,6&4&2078888403271680000&$L,1^{3},2^{4}$\\
67&A2&2,7&4&2078888403271680000&$L,2^{5},3,4$\\
68&A2&2,8&16&519722100817920000&$L,1^{3},2^{2},3^{2}$\\
69&A2&2,9&8&1039444201635840000&$L,1,2^{4},3^{2}$\\
70&A2&2,10&8&1039444201635840000&$L,1^{4},2^{2},3$\\
71&A2&2,11&8&1039444201635840000&$L,1,2^{4},3^{2}$\\
72&A2&2,14&4&2078888403271680000&$L,1,2^{5},4$\\
73&A2&2,15&4&2078888403271680000&$L,1,2^{4},3^{2}$\\
74&A2&2,16&8&1039444201635840000&$L,1,2^{4},3^{2}$\\
75&A2&2,18&4&2078888403271680000&$L,1^{2},2^{4},3$\\
76&A2&2,25&64&129930525204480000&$L^{2},1^{2},2^{4}$\\
77&A2&2,26&12&692962801090560000&$L,2^{4},3^{3}$\\
78&A2&2,28&4&2078888403271680000&$L,1^{2},2^{4},3$\\
79&A2&2,30&8&1039444201635840000&$L,1^{2},2^{4},3$\\
80&A2&2,36&48&173240700272640000&$L,2^{4},3^{3}$\\
81&A2&2,37&48&173240700272640000&$L,2^{4},3^{3}$\\
82&A2&2,38&12&692962801090560000&$L,1^{3},2^{4}$\\
83&A2&2,40&288&28873450045440000&$L^{2},2^{6}$\\
84&A2&2,41&12&692962801090560000&$L,1^{4},2^{2},3$\\
85&A2&2,42&864&9624483348480000&$L,1^{4},2^{2},6$\\
86&A2&2,44&24&346481400545280000&$L,1^{2},2^{4},3$\\
87&A2&2,46&72&115493800181760000&$L,2^{6},5$\\
88&A2&2,47&288&28873450045440000&$L,2^{6},6$\\
89&A2&2,50&24&346481400545280000&$L,1^{3},2^{3},4$\\
90&A2&2,53&288&28873450045440000&$L,2^{6},5$\\
91&A2&2,56&648&12832644464640000&$L^{2},2^{6}$\\
92&A2&2,58&1728&4812241674240000&$L^{2},1^{4},2^{2}$\\
93&A2&2,59&72&115493800181760000&$L,1^{3},2^{4}$\\
94&A2&2,60&72&115493800181760000&$L,1^{3},2^{4}$\\
95&A2&2,61&72&115493800181760000&$L,1^{3},2^{4}$\\
96&A2&2,62&648&12832644464640000&$L^{2},2^{6}$\\
97&A2&2,63&288&28873450045440000&$L^{2},1^{4},2^{2}$\\
98&A2&2,64&432&19248966696960000&$L^{2},1^{4},2^{2}$\\
99&A2&3,6&64&129930525204480000&$L,1^{2},2^{4},3$\\
100&A2&3,7&24&346481400545280000&$L,1^{4},2^{2},3$\\
101&A2&3,9&192&43310175068160000&$L,2^{4},3^{3}$\\
102&A2&3,10&240&34648140054528000&$L^{2},1^{4},2^{2}$\\
103&A2&3,11&48&173240700272640000&$L,1^{3},2^{4}$\\
104&A2&3,12&2304&3609181255680000&$L,2^{4},3^{3}$\\
105&A2&3,13&15840&524971819008000&$L^{3},1^{3},2^{2}$\\
106&A2&3,14&4320&1924896669696000&$L^{2},2^{6}$\\
107&A1&0,1&483840&17186577408000&$L^{4},1^{4}$\\
108&A1&0,2&46080&180459062784000&$L^{4},1^{4}$\\
109&A1&0,4&80640&103119464448000&$L^{4},1^{4}$\\
110&A1&0,5&2688&3093583933440000&$L^{2},1^{4},3^{2}$\\
111&A1&0,7&55296&150382552320000&$L^{4},1^{4}$\\
112&A1&0,8&3072&2706885941760000&$L^{2},1^{4},3^{2}$\\
113&A1&0,10&7920&1049943638016000&$L^{3},1^{3},2^{2}$\\
114&A1&0,11&760320&10936912896000&$L^{4},1^{4}$\\
115&A1&1,1&2304&3609181255680000&$L^{2},1^{6}$\\
116&A1&1,5&192&43310175068160000&$L,1^{6},3$\\
117&A1&1,10&96&86620350136320000&$L,1^{6},3$\\
118&A1&1,13&2880&2887345004544000&$L^{2},1^{6}$\\
119&A1&1,23&48&173240700272640000&$L,1^{5},2^{2}$\\
120&A1&1,24&24&346481400545280000&$L,1^{5},2^{2}$\\
121&A1&2,1&2304&3609181255680000&$L^{2},1^{6}$\\
122&A1&2,4&1152&7218362511360000&$L^{2},1^{6}$\\
123&A1&2,8&256&32482631301120000&$L,1^{6},3$\\
124&A1&2,10&2304&3609181255680000&$L^{2},1^{6}$\\
125&A1&2,12&80&103944420163584000&$L,1^{5},2^{2}$\\
126&A1&2,36&32&259861050408960000&$L,1^{5},2^{2}$\\
127&3&9192064&52254720&159134976000&$L^{8}$\\
128&622&416910284928&15925248&522161640000&$L^{8}$\\
129&622&5174966569752&35389440&234972738000&$L^{8}$\\
\end{longtable}
}

\subsection{Weyl-vector neighbours for all rooted Niemeier lattices}

\begin{statement}[B.6 (Refuted, with complete replacement table)]
The proposed rule that the Weyl \(d\)-neighbour has Coxeter number
\(h/d\), for every \(d\mid h\), is false. Of the 96 pairs consisting
of a rooted Niemeier lattice and a divisor of its Coxeter number,
exactly the following eleven violate the rule:
\begin{center}\small
\begin{tabular}{lrrll}
\toprule
Original root system&\(h\)&\(d\)&Neighbour root system&\(h'\ne h/d\)\\\midrule
\(D_{12}^2\)&22&2&\(D_6^4\)&\(10\ne11\)\\
\(D_8^3\)&14&2&\(D_4^6\)&\(6\ne7\)\\
\(E_8^3\)&30&2&\(D_8^3\)&\(14\ne15\)\\
\(E_8^3\)&30&3&\(A_8^3\)&\(9\ne10\)\\
\(E_8^3\)&30&5&\(A_4^6\)&\(5\ne6\)\\
\(D_6^4\)&10&2&\(A_3^8\)&\(4\ne5\)\\
\(E_6^4\)&12&3&\(A_2^{12}\)&\(3\ne4\)\\
\(D_4^6\)&6&2&\(A_1^{24}\)&\(2\ne3\)\\
\(D_{10}E_7^2\)&18&2&\(A_7^2D_5^2\)&\(8\ne9\)\\
\(D_{16}E_8\)&30&2&\(D_8^3\)&\(14\ne15\)\\
\(D_{24}\)&46&2&\(D_{12}^2\)&\(22\ne23\)\\\bottomrule
\end{tabular}
\end{center}
The complete list of all 96 neighbours, including root systems and
root counts, is \path{records/B_all_neighbours.csv}. Every
\(d=h\) neighbour is Leech.
\end{statement}

\begin{proof}
We specify the neighbour construction for arbitrary \(N\), so the
counterexamples are not an ambiguity of neighbour conventions. Choose
a positive system, let \(\rho\) be half its root sum, and let \(\alpha\)
be any simple root. For a rooted Niemeier lattice the components have
the same Coxeter number \(h\), \(\rho\in N\),
\(\rho^2=2h(h+1)\), and \((\rho,\alpha)=1\). These assertions are
respectively Lemmas 2.3, 2.6, and 2.4 of \citet{B_Borcherds1985};
in \citet{B_ChenevierLannes2019} they occur in Chapter III,
Propositions 4.2.1 and 4.2.5, with the common Coxeter number supplied
by the Niemeier root-system classification.
For \(d\mid h\) put
\[
 K_d=\{x\in N:(\rho,x)\in d\Z\},\qquad
 w_d=\rho/d-(h/d)\alpha,\qquad N_d=K_d+\Z w_d.
\]
The square of \(w_d\) is \(4h^2/d^2\), and \((w_d,K_d)\subset\Z\).
The vector \(dw_d=\rho-h\alpha\) lies in \(K_d\). Its order modulo
\(K_d\) is exactly \(d\): pairing a putative smaller multiple with
\(\alpha\), in the unimodular lattice \(N\), would require \(d\)
to divide that multiple because
\((\rho-h\alpha,\alpha)=1-2h\) is prime to \(d\).
Thus \(N_d\) is even and unimodular. Replacing \(\alpha\) by another
simple root changes \(w_d\) by an element of \(K_d\), so this
construction is independent of that choice. A Weyl change of positive
system gives an isometric neighbour.

For reproducibility the original integral Gram matrices for all 23
rooted lattices were retrieved from the catalogue of
\citet{B_NebeSloaneCatalogue}. We did not accept the labels of those
matrices as a certificate. The program \path{B_all_neighbours.py}
checks that each is positive definite, even and of determinant one,
enumerates all norm-two vectors, and verifies the stated root system.
Half the sum of the lexicographically positive roots gives an integral
\(\rho\), and the equalities above are independently checked in every
case. A root of pairing one supplies \(\alpha\). The kernel basis is
constructed by an elementary unimodular transformation of its pairing
row, followed by the single congruence condition modulo \(d\). An
exact rational span gives \(N_d\); evenness and determinant one are
checked again.

Every norm-two vector of each of the 96 resulting lattices is then
enumerated by the integer procedure already proved complete. The
connected components of the graph of nonorthogonal roots are
classified by their exact rank and number of roots. For an irreducible
simply laced positive definite root system these distinguish
\(A_n,D_n,E_6,E_7,E_8\), with the harmless equality \(A_3=D_3\).
The original and neighbour Gram matrices, Weyl vector, divisor,
and complete roots are retained for every pair. This is a finite
exhaustion of all divisors for all 23 possibilities, and yields the
table. The error in the proposed rule is the unwarranted assertion
that the quotient of a Coxeter number controls the roots surviving
or appearing in a smaller-divisor neighbour; the example
\(E_8^3,d=2\) already disproves it. Indeed there is not even an
allowed rank-24 Niemeier root system of common Coxeter number 15.

The final assertion also has a uniform proof: Chapter III,
Theorem 4.2.10(c), of \citet{B_ChenevierLannes2019} says that for an
even unimodular positive definite rank-24 lattice with roots and
Coxeter number \(h\), its Weyl \(h\)-neighbour is Leech.
Proposition 4.2.11 identifies this neighbour with the holy
construction. Our formula uses their correction vector of pairing
one and hence is precisely that neighbour. All 23 inputs satisfy
those hypotheses; the computation verifies the conclusion separately.
\end{proof}

\subsection{The Wilson cubic and an integral Niemeier embedding}

Write \(\mathbb O\) for the real division octonions with orthonormal
basis \(1,i_t\), \(t\in\Z/7\Z\), and multiplication
\(i_ti_{t+1}=i_{t+3}\), cyclically in every such triple. Put
\[
 s=(-1+i_0+\cdots+i_6)/2,
 \qquad L=D_8+\Z s,\qquad M=L\bar s,\quad H=Ls.
\]
Here \(L\) has 240 vectors of square two: the vectors
\(\pm e_i\pm e_j\), and the half-vectors with an odd number of minus
signs. Wilson's lattice is
\begin{equation}\label{eq:B-Wilson}
 \mathcal W=\{(x,y,z)\in L^3:
    x+y,y+z,z+x\in M,\ x+y+z\in H\},
\end{equation}
with squared length \((|x|^2+|y|^2+|z|^2)/2\). This is the
normalization in \citet{B_Wilson2009}. The cubic is
\(F(x,y,z)=\operatorname{Re}((xy)z)\); its polarized trilinear
expression is unambiguous because the octonion associator has real
part zero.

\begin{statement}[B.7 (Proved): the full Wilson cubic stabilizer]
The full stabilizer of \(F\) in \(\operatorname{Aut}\mathcal W\)
is
\[
 C_2^2\rtimes C_3\ \times\ (C_7\rtimes C_3)
 \ \cong\ A_4\times(C_7\rtimes C_3),
 \qquad\text{of order }252.
\]
Generators are the even sign changes of the three octonion
coordinates, their cyclic permutation, and the simultaneous Fano
permutations \(i_t\mapsto i_{t+1}\) and \(i_t\mapsto i_{2t}\).
Thus the subgroup exhibited in the workbench is the full stabilizer.
\end{statement}

\begin{proof}
First we reduce the continuous symmetry question to an explicitly
bounded finite calculation. The singular locus of \(F\) is the union
of the three coordinate octonion spaces. In fact its three derivatives
are, after octonion conjugation, the products of the other two
coordinates. The division property implies that at most one
coordinate is nonzero. A linear isometry preserving \(F\) must
therefore permute these three eight-dimensional subspaces.

The projections of \(\mathcal W\) to each block equal \(L\).
The projection of its intersection with a sum of two blocks to either
one of those blocks equals \(M\). To see this, \(M,H\subset L\)
contain \(2L\), their images modulo \(2L\) are complementary
four-dimensional subspaces, and the two-block section has
\(x,y\in M\), \(x+y\in2L\). These identities follow either directly
from the multiplication table or from the exact bases in
\path{B_wilson_stabilizer.py}. Consequently each block of a
lattice symmetry belongs to \(G=\operatorname{Aut}(L)_M\).

We determine this full projection group, rather than assuming that
familiar octonion maps generate it. The 240 roots of \(L\) are
explicit above. Choose the positive roots lexicographically and their
half sum; the eight roots pairing one with this half sum are simple.
Their root reflections generate the Weyl group of \(E_8\), which is
all \(\operatorname{Aut}(L)\), because the \(E_8\) Dynkin diagram
has no automorphisms. The exact permutation action on the roots gives
order \(696729600\). Its orbit on the subspace \(M/2L\) has size
270. The seven right multiplications \(R_{i_t}\), and the six maps
\[
 \frac12 R_{1+i_t}R_{1-i_0},\qquad t=1,\ldots,6,
\]
all preserve \(M\), have determinant \(+1\), and generate a group
of order \(2580480=696729600/270\). Orbit--stabilizer proves that
these generate the full \(G\). These checks, including the entire
finite subspace orbit, are \path{B_E8_projection.py}.

Consider now an ordered-block cubic symmetry \((A,B,C)\), with
\(A\in G\), and set \(D=\bar{\phantom{x}}\,C\,\bar{\phantom{x}}\),
where the bars denote the linear conjugation operator. The cubic
identity is equivalent to
\begin{equation}\label{eq:B-triality}
 (Ax)(By)=D(xy)\quad(x,y\in\mathbb O).
\end{equation}
For a fixed \(A\) there are at most two orthogonal solutions
\((B,C)\). Indeed the ratio of two solutions has first component
identity. Equation \eqref{eq:B-triality} with \(x=1\), then \(y=1\),
shows that its remaining maps are right multiplication by one unit
\(b\), and that \(x(yb)=(xy)b\) for all \(x,y\). The nucleus of
\(\mathbb O\) is \(\R\), hence \(b=\pm1\). In the supplied
multiplication table this last statement can also be verified by the
64 linear equations in the eight coordinates of \(b\); their
solution space is exactly \(\R\cdot1\).

For each of the thirteen generators of \(G\) we solve
\eqref{eq:B-triality} exactly. If \(a=A1\), then \(D=L_aB\), so
\((Ax)(By)=aB(xy)\) is a system of 512 rational linear equations in
the 64 entries of \(B\). In each case its nullspace is one dimensional
and admits a rational orthogonal normalization. The matrices are
recorded in \path{B_triality_lifts.json}. Adjoin the kernel
\((I,-I,-I)\). Their faithful permutation action on the roots in
the first and second blocks has order \(5160960=2|G|\).
This proves that we have exhausted all possible first projections
and their two lifts.

It remains to impose the lattice conditions on the other blocks.
The orbit of \(L\) under the third components of this lift group
has size 15, giving a stabilizer of order \(344064\). Its Schreier
generators are obtained by exact rational lattice Hermite normal
forms. On this subgroup all three components preserve \(L\).
Modulo \(2L^3\), the glue of \(\mathcal W\) is the twelve-dimensional
binary code
\[
 \{(m_1+n,m_2+n,m_1+m_2+n):
        m_1,m_2\in M/2L,\ n\in H/2L\}.
\]
The orbit of this code has size 4096. The ordered-block stabilizer
therefore has order
\(5160960/(15\cdot4096)=84\). Exact Schreier generators are retained
in \path{B_wilson_stabilizer.json}; their generated order is
independently checked in their permutation action. The code orbit
is exhaustive: a queue starts at the displayed code and closes under
every generator, comparing reduced binary echelon bases exactly.
Thus the order is an upper bound as well as a lower bound.

No odd permutation of blocks can occur. After a cyclic permutation
and composition with one of the even triality lifts having the same
first block, such a symmetry would have first block \(I\) and
interchange the other two. Its identity would force
\(x(yb)=(yx)b\) for all \(x,y\), for a unit \(b\). This is
impossible. For completeness the same exact linear calculation gives
a zero-dimensional solution space in the eight coordinates of \(b\),
whereas the non-reversed equation above gives dimension one. This is
a finite verification of the obstruction over \(\Q\), and hence
over \(\R\), since the coefficients are rational.

Finally the listed even sign changes, coordinate cycle and two Fano
permutations preserve \eqref{eq:B-Wilson} and the cubic. The first
two operations generate \(A_4\); the Fano permutations fix \(s\),
preserve the multiplication table, and generate \(C_7\rtimes C_3\).
The two groups commute and have trivial intersection. They give
\(4\cdot3\cdot21=252\) elements, precisely three times the
ordered-block upper bound 84. This proves both fullness and the
claimed isomorphism type.
\end{proof}

\begin{statement}[B.8 (Proved): a Niemeier embedding with integral twice cubic]
The Niemeier lattice \(E_8^3\), with its usual norm-two roots, has the
explicit isometric embedding \(L^3\subset\mathbb O^3\), where
\(L=D_8+\Z s\) as above and the Euclidean norm is
\(|x|^2+|y|^2+|z|^2\). Its cubic value group, meaning the additive
subgroup generated by all \(F(x,y,z)\), is
\[
 \langle F(L^3)\rangle_{\Z}=\tfrac12\Z.
\]
In particular \(2F\) is integer valued on this embedding. This
arithmetic property depends on the embedding: rotating the real and
\(i_0\) axes of the first block through \(\pi/4\) gives another
isometric embedding whose cubic value group is not \(\tfrac12\Z\).
More precisely, the second value group is
\(\tfrac12\Z+\tfrac{\sqrt2}{4}\Z\).
\end{statement}

\begin{proof}
Here the ambient metric is the ordinary octonion sum norm, the
coefficient of \(-\lambda\) in the Albert determinant
\(\lambda^3-\lambda(|x|^2+|y|^2+|z|^2)+2F(x,y,z)\).
Wilson's separately specified half-sum norm in B.7 is identified
with this metric by \((x,y,z)\mapsto(x,y,z)/\sqrt2\).
That change multiplies the cubic by \(1/(2\sqrt2)\) and leaves
its stabilizer unchanged; B.8 uses the ordinary metric throughout.
The lattice \(L\) is even unimodular: \(D_8\) has determinant four,
\(s^2=2\), \((s,D_8)\subset\Z\), and adjoining \(s\) has index
two. Thus its orthogonal cube is the stated Niemeier lattice.
The real trilinear expression is cyclic in its three arguments.
If all arguments belong to \(D_8\), its value is integral; if one
argument is \(s\), it lies in \(\tfrac12\Z\). If two arguments
are \(s\), cyclicity makes them adjacent and the coordinates of
\(s^2\) are half integers, so the same conclusion follows. Finally
\(\operatorname{Re}(s^3)=5/2\). Trilinearity in a generating set
therefore proves \(F(L^3)\subset\tfrac12\Z\). On the other hand
\((s,s,s)\) gives \(5/2\), and \((2,2,2)\) gives 8, so their
integer combinations generate \(\tfrac12\Z\).

Let \(T\) act orthogonally on the first octonion block by
\(T1=(1+i_0)/\sqrt2\), \(Ti_0=(-1+i_0)/\sqrt2\), fixing the
remaining six coordinates. The second embedding is
\((x,y,z)\mapsto(Tx,y,z)\). Its value at \((2,2,2)\) is
\(4\sqrt2\), which is not in \(\tfrac12\Z\). Thus the arithmetic
criterion offered in B.4 is satisfied by an explicit embedding, but
does not make the value group an invariant of the abstract Niemeier
lattice. No assertion of uniqueness or canonicality of that embedding
is implied by the word ``distinguished.''

We also compute the entire second group, rather than only an
irrational witness. Write \(e_0=1,e_{t+1}=i_t\), and take the
\(E_8\) basis
\(s,d_1,\ldots,d_7\), where
\(d_j=e_j-e_{j+1}\) for \(1\le j\le6\), and
\(d_7=e_6+e_7\). Its Gram determinant is one. Trilinearity reduces
the value group to the 512 triples of these basis vectors.
In \path{B_cubic_embeddings.py} every value is computed exactly
as a pair of rational coefficients of \(1,\sqrt2\); all have
first coefficient in \(\tfrac12\Z\) and second in
\(\tfrac14\Z\). Conversely the three exact values
\[
 F(Ts,s,d_4)=\tfrac12,\qquad
 F(Ts,d_7,s)=\tfrac{\sqrt2}{2},\qquad
 F(Ts,s,s)=\tfrac32+\tfrac{3\sqrt2}{4}
\]
generate \(\tfrac12\Z+\tfrac{\sqrt2}{4}\Z\). The complete
coefficient list and its two-dimensional integer Hermite reduction
are retained in \path{B_cubic_embeddings.json}. Thus both embeddings
and both distinct value groups have exact finite certificates.
\end{proof}

\paragraph{Literature scope.}
Borcherds' construction and the Weyl-neighbour results of Chenevier and
Lannes supply the rank-24 background and the endpoint \(d=h\); they
do not supply the eleven-exception table. Wilson's 2009 construction
supplies \eqref{eq:B-Wilson}, and his account of Conway's group supplies
large generating groups preserving that lattice. A generator of
\(\mathrm{Co}_0\) need not preserve the cubic, which is why the
separate upper-bound argument in B.7 is required. Searches for the
related Albert-algebra interpretation also locate
\citet[Corollary 2.3 and Section 4]{B_Rios2018}. Its equation (8)
uses the quadratic map \(v\mapsto vv^*\) into the full Albert
algebra. That construction does not identify the fixed linear
off-diagonal coordinates with a cubic-preserving action on
\(\mathbb O^3\); in particular \(v\) and \(-v\) have the same
quadratic image, whereas our nonzero cubic changes sign. We
therefore do not infer a cubic stabilizer from that interpretation.
Searches for the specific determinant-81 intersection, its shell counts 23058 and
1849032, and its group orders, together with searches for isotropic
planes modulo three, did not identify a published occurrence of this
configuration or a full plane-orbit table. We therefore make no
priority claim and do not label the configuration new. The exact
searches, sources opened, and theorem locations are retained in the
Part B research record. The full 129-orbit table in B.9 is an
independent exact classification; the failure to locate an earlier
table is not a claim of priority. The vector-residue table in
\citet[Proposition 4.2 and Table 4.2]{B_MartinetLatticesMod} and the
explicit vector transporters documented by \citet{B_mmgroup2026}
are useful comparisons, but neither supplies the plane fusion or
the 129 full stabilizers proved here. The known triangle stabilizer
in \citet[Section 5]{B_Fairbairn2010} is independently recovered by
our complete Gram-graph coset test. The Albert/Niemeier connection
itself is present in the work of \citet{B_ElkiesGross2001}; B.8
addresses the particular Euclidean embedding and cubic normalization
asked for here, and its elementary value-group calculation makes no
claim to introduce that broader connection.

The requested Conway--Sloane comparison is more specific than a
general reference to their book. Their Chapter 16 tables record the
rooted Niemeier systems and glue; the opened Chapter 17 material
classifies odd unimodular lattices through their two even neighbours,
an index-two construction. For Chapter 24 we opened the complete
original article \citep{B_ConwaySloane1982Holy}, whose Figure 3
includes both \(A_5^4D_4\) and \(A_2^{12}\). Its unnumbered holy
construction on p.~275, proved on pp.~281--283, constructs a Leech
lattice from each rooted Niemeier lattice, with the norm-two root
normalization used here. This is known background for B.1 and B.3.
The opened material does not state the common determinant-81
intersection, its six Leech and two rooted extensions, or the
129 plane orbits. We distinguish that precise search outcome from
an assertion that the configuration is absent from all literature.
